% ECAC Confidentiality & Key Rotation (M9)
% Figure 6: For Section 3.4 (Confidentiality)

\begin{tikzpicture}[
    % Styles
    databox/.style={
        rectangle, rounded corners=4pt,
        minimum width=3cm,
        font=\tiny\sffamily, align=left,
        draw=black!50, fill=white, line width=0.5pt,
        inner sep=5pt
    },
    opbox/.style={
        rectangle, rounded corners=4pt,
        minimum width=3.5cm,
        font=\tiny\sffamily, align=left,
        draw=#1!60!black, fill=#1!15, line width=0.5pt,
        inner sep=5pt
    },
    process/.style={
        rectangle, rounded corners=4pt,
        minimum width=2.2cm, minimum height=0.8cm,
        font=\scriptsize\sffamily, align=center,
        draw=#1!60!black, fill=#1!15, line width=0.6pt
    },
    decision/.style={
        diamond, aspect=2,
        minimum width=2cm, minimum height=1cm,
        font=\scriptsize\sffamily, align=center,
        draw=#1!60!black, fill=#1!15, line width=0.6pt
    },
    arrow/.style={
        -{Stealth[length=2mm, width=1.5mm]},
        line width=0.7pt, draw=black!60
    },
    timeline/.style={
        line width=1.5pt, draw=black!40
    },
    event/.style={
        circle, minimum size=0.5cm,
        font=\tiny\sffamily\bfseries,
        draw=#1!70!black, fill=#1!30, line width=0.5pt
    },
    zone/.style={
        rectangle, rounded corners=8pt,
        draw=black!40, fill=black!3, line width=0.8pt,
        inner sep=8pt
    },
    zonetitle/.style={
        font=\small\sffamily\bfseries, text=black!70
    }
]

% Colors
\definecolor{plain}{HTML}{3498DB}
\definecolor{cipher}{HTML}{9B59B6}
\definecolor{key}{HTML}{F39C12}
\definecolor{grant}{HTML}{27AE60}
\definecolor{revoke}{HTML}{E74C3C}
\definecolor{redact}{HTML}{95A5A6}

% === TOP LEFT: PER-TAG ENCRYPTION ===
\node[zonetitle] at (3.5, 9) {PER-TAG ENCRYPTION};

% Plaintext op
\node[databox, fill=plain!10] (plainop) at (1.5, 7.5) {
    \textbf{Data op (plain):}\\
    key: ``doc:1:body''\\
    value: ``secret text''\\
    tag: ``classified''
};

% Arrow
\node[font=\scriptsize\sffamily, text=cipher!70!black] at (3.5, 7.5) {encrypt\\$(K_{tag})$};
\draw[arrow, color=cipher!70!black, line width=1pt] (2.8, 7.5) -- (4.2, 7.5);

% Ciphertext op
\node[databox, fill=cipher!10] (cipherop) at (5.5, 7.5) {
    \textbf{Data op (encrypted):}\\
    key: ``doc:1:body''\\
    value: $\langle$ciphertext$\rangle$\\
    tag: ``classified''
};

% Crypto details
\node[font=\tiny\sffamily, text=black!50, align=left] at (3.5, 6.2) {
    Cipher: XChaCha20-Poly1305\\
    Overhead: +40 bytes\\
    AAD: blake3(obj $\|$ field $\|$ op\_id)
};

% Key lookup
\node[font=\tiny\sffamily, text=key!70!black] at (3.5, 5.6) {
    Key lookup: tag $\rightarrow$ (version, $K_{sym}$)
};

% === TOP RIGHT: KEY DISTRIBUTION ===
\node[zonetitle] at (11, 9) {KEY DISTRIBUTION \& ROTATION};

% KeyRotate op
\node[opbox=key] (keyrot) at (11, 7.8) {
    \textbf{KeyRotate} \{\\
    \quad tag: ``classified''\\
    \quad new\_version: 2\\
    \quad new\_key: $E(K_2, \text{admin.pk})$\\
    \}
};

% KeyGrant op
\node[opbox=grant] (keygrant) at (11, 5.8) {
    \textbf{KeyGrant} \{\\
    \quad subject: Alice.pk\\
    \quad tag: ``classified''\\
    \quad key\_version: 2\\
    \quad encrypted\_key: $E(K_2, \text{Alice.pk})$\\
    \}
};

% Arrow between
\draw[arrow, line width=1pt] (keyrot) -- node[right, font=\tiny\sffamily] {then} (keygrant);

% === MIDDLE: FORWARD SECRECY TIMELINE ===
\node[zonetitle] at (7, 4.2) {FORWARD SECRECY TIMELINE};

% Timeline
\draw[timeline] (1, 3) -- (13, 3);
\draw[line width=0.8pt, draw=black!40] (1, 2.9) -- (1, 3.1);
\draw[line width=0.8pt, draw=black!40] (13, 2.9) -- (13, 3.1);
\node[font=\tiny\sffamily, text=black!50] at (13.3, 3) {$t$};

% Events on timeline
\node[event=grant] (grantev) at (3, 3) {};
\node[font=\tiny\sffamily, anchor=south] at (3, 3.3) {Grant(Bob)};

\node[event=revoke] (revokeev) at (6, 3) {};
\node[font=\tiny\sffamily, anchor=south] at (6, 3.3) {Revoke(Bob)};

\node[event=key] (rotev) at (9, 3) {};
\node[font=\tiny\sffamily, anchor=south] at (9, 3.3) {KeyRotate};
\node[font=\tiny\sffamily, anchor=south] at (9, 3.7) {($K_1 \rightarrow K_2$)};

% Key regions
\draw[decorate, decoration={brace, amplitude=4pt}, line width=0.5pt]
    (1, 2.6) -- (9, 2.6) node[midway, below=4pt, font=\tiny\sffamily] {encrypted with $K_1$};
\draw[decorate, decoration={brace, amplitude=4pt}, line width=0.5pt]
    (9, 2.6) -- (13, 2.6) node[midway, below=4pt, font=\tiny\sffamily] {encrypted with $K_2$};

% Bob's access
\draw[line width=2pt, color=grant!60!black] (3, 2.2) -- (6, 2.2);
\node[font=\tiny\sffamily, text=grant!70!black, anchor=east] at (2.8, 2.2) {Bob reads};

\draw[line width=2pt, color=revoke!60!black, dashed] (6, 2.2) -- (9, 2.2);
\node[font=\tiny\sffamily, text=revoke!70!black] at (7.5, 1.8) {Bob loses access};

% Forward secrecy annotation
\node[zone, minimum width=3cm, minimum height=1cm, fill=revoke!5] at (11, 2) {};
\node[font=\tiny\sffamily, align=center, text=revoke!70!black] at (11, 2) {
    Bob cannot read\\
    new data ($K_2$)\\
    \textbf{forward secrecy}
};

% === BOTTOM: QUERY WITH REDACTION ===
\node[zonetitle] at (7, 0.5) {QUERY WITH REDACTION};

% Query process
\node[process=plain] (queryproc) at (2, -0.8) {project\_field\\for\_subject()};

% Decision
\node[decision=key] (checkepoch) at (5.5, -0.8) {KeyGrant\\epoch?};
\draw[arrow] (queryproc) -- (checkepoch);

% Outcomes
\node[process=grant] (decrypt) at (9, 0) {Decrypt\\$\rightarrow$ plaintext};
\node[process=redact] (redact) at (9, -1.6) {Return\\``$\langle$redacted$\rangle$''};

\draw[arrow, color=grant!70!black] (checkepoch.north) |- node[above, pos=0.7, font=\tiny\sffamily] {yes} (decrypt.west);
\draw[arrow, color=redact!70!black] (checkepoch.south) |- node[below, pos=0.7, font=\tiny\sffamily] {no} (redact.west);

% Output examples
\node[font=\tiny\ttfamily, text=grant!70!black, anchor=west] at (10.5, 0) {``secret text''};
\node[font=\tiny\ttfamily, text=redact!70!black, anchor=west] at (10.5, -1.6) {``<redacted>''};

% Note
\node[font=\tiny\sffamily, text=black!50, align=center] at (7, -2.8) {
    Unauthorized subjects see consistent state structure but cannot read encrypted field values
};

\end{tikzpicture}
